\section{The generator, and what it contains}\label{sec:generator}

Definition~\ref{def:generator} names an operator; this section identifies it. The
outcome is that \eqref{eq:form} holds with the \emph{complete} target on the right,
and that the whole $\omega$-dependence is one scalar factor.

\begin{theorem}[The shifted form in closed form]\label{thm:kernel}
For $0\leq\omega\leq\frac12$ and $f\in C_c^\infty(I_L)$, $F=E_Lf$,
\begin{align}
 Q_{\omega,L}[f]
 ={}&\int_0^\infty\big[2\norm F^2-2\cosh(\omega r)\re\ip F{U_rF}\big]
      \gammak(r)\dd r
   \;+\;w_0\norm f^2 \notag\\
  &+2\iint_{I_L^2}\cosh\!\big(\omega(x-y)\big)
      \cosh\!\Big(\frac{x-y}2\Big)\overline{f(x)}f(y)\dd x\dd y \notag\\
  &-2\!\!\sum_{m\log p<L}\!\!(\log p)\,p^{-m/2}\cosh(\omega m\log p)\,
      \re\ip F{U_{m\log p}F} .
 \label{eq:shifted-target}
\end{align}
At $\omega=0$ this is \eqref{eq:target}. In words: \emph{the shifted form is the
central form with every off-diagonal kernel multiplied by $\cosh(\omega(x-y))$
and every prime atom by $\cosh(\omega\log n)$, the diagonal unchanged.}
\end{theorem}

\begin{proof}
Write $\mathcal F(p)=\xi(\frac12+p)$ and $s=\frac12+p$. Logarithmic
differentiation of \eqref{eq:xi} gives
\begin{equation}\label{eq:threefactors}
 \frac{\mathcal F'}{\mathcal F}(p)
 =\underbrace{\frac1{p+\frac12}+\frac1{p-\frac12}}_{\frac12s(s-1)}
 \;\underbrace{-\tfrac12\log\pi+\tfrac12\psi\big(\tfrac14+\tfrac p2\big)}
   _{\pi^{-s/2}\Gamma(s/2)}
 \;+\;\underbrace{\frac{\zeta'}{\zeta}\big(\tfrac12+p\big)}_{\zeta(s)},
\end{equation}
and $a_\omega(p)$ of \eqref{eq:gensymbol} is the sum of \eqref{eq:threefactors}
at $p-\omega$ and at $p+\omega$. We compute the causal kernel of each factor on
$\re p=\eta>1$ and then symmetrize. For a causal convolution with kernel $c$ and
$g$ supported in $I_L$,
\begin{equation}\label{eq:symmetrize}
 \re\ip g{(c*\widetilde g)|_{I_L}}
 =\tfrac12\iint_{x\neq y}c(|x-y|)\overline{g(x)}g(y)\dd x\dd y ,
\end{equation}
since relabelling $x\leftrightarrow y$ in the conjugate term turns the region
$y<x$ into $x<y$ with $c(|x-y|)$ unchanged.

\emph{Pole factors.} The four terms $\frac1{p\mp\omega\pm\frac12}$ have simple
poles at $p=\pm\frac12\pm\omega$, all to the left of $\re p=\eta$ since
$\eta>1\geq\frac12+\omega$. Because $\frac1{p-a}$ is the Laplace transform of
$e^{ax}\mathbf1_{x>0}$, their kernel is
\[
 \big(e^{x/2}+e^{-x/2}\big)\big(e^{\omega x}+e^{-\omega x}\big)\mathbf1_{x>0}
 =4\cosh(x/2)\cosh(\omega x)\mathbf1_{x>0},
\]
and \eqref{eq:symmetrize} turns it into
$2\iint\cosh(\omega(x-y))\cosh(\frac{x-y}2)\overline f f$, the third line of
\eqref{eq:shifted-target}. At $\omega=0$ the expansion
$2\cosh\frac{x-y}2=2\cosh\frac x2\cosh\frac y2-2\sinh\frac x2\sinh\frac y2$
returns $P_L[f]$ of \eqref{eq:poles}.

\emph{Zeta factor.} On $\re p=\eta>1$ the Dirichlet series converges absolutely,
so
\[
 \frac{\zeta'}{\zeta}\big(\tfrac12+p-\omega\big)
 +\frac{\zeta'}{\zeta}\big(\tfrac12+p+\omega\big)
 =-2\sum_{n\geq2}\Lambda(n)n^{-1/2}\cosh(\omega\log n)\,n^{-p},
\]
and $n^{-p}=e^{-p\log n}$ is the Laplace symbol of the causal delay by $\log n$.
Applying \eqref{eq:symmetrize} to the atom at $r=\log n$ gives
$-2\Lambda(n)n^{-1/2}\cosh(\omega\log n)\re\ip F{U_{\log n}F}$; writing
$n=p^m$, $\Lambda(p^m)=\log p$ gives the last line of
\eqref{eq:shifted-target}. The cutoff $m\log p<L$ is automatic, because a delay
of length $\geq L$ has zero action on $I_L$.

\emph{Gamma and $\pi$ factors.} Their contribution to $a_\omega$ is
$-\log\pi+\frac12[\psi(\frac14+\frac{p-\omega}2)+\psi(\frac14+\frac{p+\omega}2)]$.
With $\psi(z)=-\gamma_E+\int_0^\infty(e^{-t}-e^{-zt})(1-e^{-t})^{-1}\dd t$
\cite[\S5.9]{DLMF} at $z=\frac14+\frac{p\mp\omega}2$ and the substitution
$t=2r$,
\begin{equation}\label{eq:digamma-kernel}
 \tfrac12\Big[\psi\big(\tfrac14+\tfrac{p-\omega}2\big)
  +\psi\big(\tfrac14+\tfrac{p+\omega}2\big)\Big]
 =-\gamma_E+\int_0^\infty\frac{2e^{-2r}}{1-e^{-2r}}\dd r
  -\int_0^\infty 2\cosh(\omega r)\gammak(r)e^{-pr}\dd r ,
\end{equation}
using $\gammak(r)=e^{-r/2}(1-e^{-2r})^{-1}$ of \eqref{eq:gamma-kernel}. The two
integrals in \eqref{eq:digamma-kernel} diverge logarithmically at the origin and
only their difference converges, both integrands being $\sim\frac1{2r}$; the
pairing of the divergent constant against $\norm f^2$ with the divergent
convolution against $\re\ip F{U_rF}$ is exactly the combination
$2\norm F^2-2\cosh(\omega r)\re\ip F{U_rF}$ appearing in the first line of
\eqref{eq:shifted-target}, which is integrable against $\gammak$ because it is
$O(r)\gammak(r)$ at the origin and $2\norm F^2\gammak(r)$ for $r\geq L$. This
fixes the first line up to an $\omega$-independent constant multiple of
$\norm f^2$, and that constant is $w_0=\psi(\frac14)-\log\pi$ by evaluating at
$\omega=0$, $p=i\tau$, where the multiplier of \eqref{eq:gamma-energy} and
\eqref{eq:multiplier} is $\re\psi(\frac14+\frac{i\tau}2)-\log\pi=b(\tau^2)+w_0$.

Each of the three contributions is $\omega$-independent on the diagonal, since
$\cosh0=1$; summing them gives \eqref{eq:shifted-target}, and at $\omega=0$
gives \eqref{eq:target}. The $\omega$-dependent part of each line is separately
convergent, because $\cosh(\omega t)-1=O(t^2)$ at the origin.
\end{proof}

\begin{remark}[The equations of the background are general]\label{rem:general}
Theorem~\ref{thm:kernel} settles a question that the background leaves implicit.
The flow identities \eqref{eq:flow}--\eqref{eq:storage} are statements about the
\emph{whole} Weil form: the archimedean energy, the contact, the poles and the
primes all appear, and all are deformed. A separate paragraph of the background
isolates the gamma and $\pi$ factors alone, $K_\omega^\Gamma$, and observes that
a positive tower realizes only that portion of the zero-shift generator; that is
a statement about a construction, not about the scope of \eqref{eq:flow}.
\end{remark}

\begin{remark}[The pole term is invisible on the imaginary axis]\label{rem:imaginary}
The causal definition on $\re p=\eta>1$ is not a technical convenience. At
$p=i\tau$,
\[
 \re\left[\frac2{p-\frac12}+\frac2{p+\frac12}\right]
 =\frac{2\cdot\frac12}{\frac14+\tau^2}+\frac{2\cdot(-\frac12)}{\frac14+\tau^2}
 =0
\]
identically in $\tau$. Reading $a_\omega$ as a Fourier multiplier on the
imaginary axis therefore recovers the archimedean term and the prime delays and
\emph{loses the contact and pole term entirely}. The pole factors reach
$Q_{\omega,L}$ only through the right half-plane realization of
Definition~\ref{def:generator}. Appendix~\ref{app:checks} verifies both halves of
this remark numerically: the causal kernel $4\cosh(x/2)\mathbf1_{x>0}$
reproduces $P_L$ to machine precision, and the same symbol has vanishing real
part on the imaginary axis.
\end{remark}
